\section{Chapitre 2 : Analyse des Besoins et Spécifications}
\addcontentsline{toc}{section}{Chapitre 2 : Analyse des Besoins et Spécifications}

La phase d'analyse des besoins constitue le fondement épistémologique et méthodologique de tout projet d'ingénierie logicielle d'envergure, car elle permet de circonscrire précisément le périmètre d'action du système à concevoir et de traduire les attentes disparates des diverses parties prenantes en spécifications techniques rigoureuses. Dans le cadre de la plateforme d'apprentissage intelligent Fraudly, cette étape s'avère particulièrement cruciale en raison de l'intrication complexe entre des modules fonctionnels hétérogènes, allant de la gestion classique de ressources pédagogiques à l'orchestration avancée d'agents d'intelligence artificielle et à l'analyse vidéo en temps réel pour la prévention de la fraude. La première étape de cette démarche analytique a consisté à identifier exhaustivement les profils d'utilisateurs qui interagiront avec le système, ainsi que la nature de leurs prérogatives respectives. La plateforme s'articule autour de trois acteurs principaux, dont les parcours d'utilisation dictent l'architecture fonctionnelle de la solution. L'étudiant représente l'utilisateur final et le bénéficiaire central du système, dont l'objectif premier est d'acquérir des compétences de manière personnalisée et d'être évalué équitablement. L'enseignant occupe le rôle d'ingénieur pédagogique et de superviseur, chargé d'alimenter la base de connaissances, de paramétrer les règles d'évaluation et de valider les rapports d'intégrité. Enfin, l'administrateur système assume la responsabilité de la maintenance technique, de la gestion des habilitations globales et de l'assurance de la stabilité de l'infrastructure technologique sous-jacente.

Concernant les spécifications fonctionnelles dévolues à l'étudiant, la plateforme doit avant tout offrir une interface immersive et intuitive de consultation des parcours d'apprentissage. L'étudiant doit pouvoir s'authentifier de manière sécurisée et accéder à un tableau de bord personnalisé récapitulant ses cours en cours, ses chapitres terminés et ses lacunes identifiées. Il est impératif que le système ne se contente pas de présenter une liste statique de ressources, mais qu'il génère dynamiquement des recommandations de modules à réviser, orchestrées par l'agent d'apprentissage adaptatif en fonction des scores obtenus lors des évaluations précédentes. L'étudiant doit pouvoir interagir avec le matériel pédagogique (fichiers au format PDF, séquences vidéo, documents textuels) et bénéficier de l'assistance en temps réel d'un agent tuteur virtuel. Cet agent conversationnel, incarné par une interface de discussion intégrée, doit être capable de répondre aux interrogations conceptuelles de l'étudiant en restreignant ses réponses au strict contexte sémantique du cours, prévenant ainsi les hallucinations factuelles communément observées dans les modèles de langage génériques. Lorsque l'étudiant aborde une phase d'évaluation, le système doit basculer dans un environnement hautement sécurisé. Le déclenchement d'un examen implique obligatoirement l'initialisation d'une session de télésurveillance, requérant le consentement explicite de l'étudiant pour l'activation de sa webcam et la capture périodique de son écran. Durant l'épreuve, l'étudiant doit naviguer à travers une série de questions qui ont été potentiellement générées sur mesure par l'intelligence artificielle, comprenant des questions à choix multiples, des assertions de type vrai ou faux, ainsi que des questions ouvertes nécessitant une rédaction argumentée. Une fois l'évaluation soumise, l'étudiant s'attend à recevoir une correction détaillée, argumentée et accompagnée de recommandations pour consolider ses acquis défaillants.

Du point de vue de l'enseignant, la plateforme Fraudly doit se comporter comme un outil d'assistance à la scénarisation pédagogique qui décharge considérablement le corps professoral des tâches répétitives et chronophages. L'enseignant doit disposer de fonctionnalités complètes pour la création et l'arborescence des cours, lui permettant de structurer ses enseignements en chapitres logiques et de téléverser des ressources multimédias diversifiées. Une exigence fonctionnelle majeure réside dans l'intégration d'un mécanisme de reconnaissance optique de caractères automatique lors du téléversement de documents non textuels, afin que le contenu extrait puisse être ingéré et vectorisé par la base de données sémantique de l'intelligence artificielle. En matière d'évaluation, l'enseignant doit pouvoir solliciter l'agent de génération d'évaluations en spécifiant simplement des paramètres de haut niveau tels que le chapitre visé, le nombre de questions souhaité, le niveau de difficulté globale et les types de questions désirés. Le système doit alors proposer un projet d'examen que l'enseignant peut réviser et valider avant publication. Par la suite, le processus de correction des questions à réponse ouverte doit être délégué de manière asynchrone à un agent correcteur, qui évaluera les propositions des étudiants en s'appuyant sur des rubriques de notation explicites fournies préalablement par l'enseignant. Après la clôture d'une session d'examen, l'enseignant doit avoir accès à un tableau de bord d'intégrité complet généré par le module de télésurveillance. Ce rapport ne doit pas se limiter à une énumération brute d'incidents techniques, mais doit présenter un score de fiabilité assorti d'un compte rendu analytique distinguant les anomalies bénignes des tentatives de fraude caractérisées, avec des liens directs vers les preuves visuelles ou comportementales incriminantes (captures d'écran, enregistrements vidéo des événements critiques). L'enseignant conserve l'autorité souveraine pour invalider ou confirmer la décision suggérée par le système automatisé.

Au-delà de ces spécifications fonctionnelles exhaustives, l'ingénierie de la plateforme Fraudly est soumise à un cahier des charges non fonctionnel des plus stricts, dicté par la sensibilité des données manipulées et la criticité du contexte académique. La première de ces exigences concerne la haute disponibilité et la scalabilité horizontale du système. Les sessions d'évaluation se déroulant typiquement à des créneaux horaires fixes et impliquant simultanément des centaines, voire des milliers d'étudiants, l'architecture technologique doit garantir une tolérance absolue aux pics de charge sans dégradation perceptible de la latence, particulièrement en ce qui concerne les appels aux modèles d'intelligence artificielle et l'analyse vidéo en continu. La sécurité de l'information constitue une autre préoccupation prééminente. Le système doit implémenter une politique de contrôle d'accès rigoureuse basée sur les rôles, sécuriser les échanges de données via des protocoles cryptographiques éprouvés, et assurer la pseudonymisation des données personnelles afin de se conformer scrupuleusement aux législations en vigueur sur la protection de la vie privée. L'intégrité des flux vidéo collectés lors de la télésurveillance exige des mesures de protection contre la falsification, et ces données sensibles doivent être purgées automatiquement à l'issue des délais légaux de conservation ou des périodes de recours académique. Enfin, la modularité et la maintenabilité s'imposent comme des directives architecturales fondamentales. La volatilité intrinsèque du domaine de l'intelligence artificielle générative requiert que le couplage entre la logique applicative métier et les implémentations spécifiques des grands modèles de langage soit le plus lâche possible, permettant ainsi l'interchangeabilité aisée des fournisseurs cognitifs au gré des évolutions technologiques et des contraintes économiques futures.
